\documentclass[11pt]{article}
\usepackage[a4paper,margin=2.5cm]{geometry}
\usepackage{graphicx}
\usepackage{booktabs}
\usepackage{longtable}
\usepackage{hyperref}
\usepackage{xcolor}
\usepackage{listings}
\usepackage{parskip}
\usepackage{titlesec}
\usepackage{caption}

\hypersetup{colorlinks=true, linkcolor=blue, urlcolor=blue, citecolor=blue}
\titleformat{\section}{\Large\bfseries}{\thesection}{1em}{}
\titleformat{\subsection}{\large\bfseries}{\thesubsection}{1em}{}

\title{\textbf{MIMaaS: Extending a Machine-Intelligence-as-a-Service Platform\\
from Microcontrollers to Edge NPUs}\\[0.3em]
\large Internship Final Report}
\author{Maxime \\ AnkiLab}
\date{August 2026}

\begin{document}
\maketitle

\begin{abstract}
This report summarizes the work carried out during an internship at AnkiLab on the
\emph{MIMaaS} (Machine Intelligence with Microcontroller/NPU as a Service) platform.
The project started from an existing service that benchmarks TensorFlow Lite models
on microcontrollers (MCUs) over Zephyr RTOS, and extended it to support the
Hailo-10H edge NPU accelerator on a Raspberry Pi 5. The work covered restructuring
of the existing MCU codebase, fixing Zephyr/TFLite-Micro compatibility issues,
building a new hardware backend for Hailo, and running a systematic compatibility
and performance study of 30+ neural network architectures on the Hailo Dataflow
Compiler (DFC) 5.3.0 toolchain. Of the models tested, 14 pre-compiled reference
networks passed with 100\% success, and of 24 additional TensorFlow/Keras
architectures (17 common CNN backbones plus 7 small, purpose-built models each
probing a single architectural feature) converted to TFLite, 16 compiled and
ran successfully while 8 failed for seven distinct, well-characterized root
causes. This report presents the
architecture, methodology, quantitative results, and a discussion of the
recurring failure patterns observed, together with recommendations for future work.
\end{abstract}

\tableofcontents
\newpage

\section{Introduction and Motivation}

Edge inference of deep neural networks is increasingly deployed across a wide
spectrum of hardware, from ultra-low-power microcontrollers (MCUs) to dedicated
neural processing units (NPUs). Developers and researchers who want to evaluate
whether a given model is deployable on a given piece of hardware are typically
forced to install vendor-specific toolchains, understand board-specific
flashing procedures, and manually instrument power/latency measurements. This is
a significant barrier, especially for teams that want to compare multiple
hardware targets before committing to one.

\emph{MIMaaS} (Machine Intelligence with Microcontroller-as-a-Service) is a
platform developed at AnkiLab that removes this barrier: a user submits a
\texttt{.tflite} model through a REST API, and the server automatically builds
firmware (or a hardware-specific binary), flashes it to physical hardware,
executes inference, and returns hardware-measured metrics (inference time, RAM,
ROM/flash usage, and, where instrumented, power and energy per inference).

Prior to this internship, MIMaaS supported only microcontroller-class boards
(Nordic nRF52/nRF53 series, STM32 Nucleo, ESP32-C6) driven by Zephyr RTOS and
TensorFlow Lite for Microcontrollers (TFLite Micro). The motivation for this
internship was twofold:

\begin{enumerate}
  \item \textbf{Consolidate and repair} the existing MCU pipeline, which had
    accumulated technical debt (a flat, monolithic codebase) and had broken
    after a Zephyr version migration.
  \item \textbf{Extend MIMaaS to a new hardware class}: edge NPUs, specifically
    the Hailo-10H accelerator connected to a Raspberry Pi 5, to evaluate whether
    the same request/response abstraction (``upload a model, get back
    performance numbers'') generalizes beyond microcontrollers to
    dataflow-compiled NPU accelerators.
\end{enumerate}

\section{Problem Statement}

Two concrete problems were addressed:

\textbf{Problem 1 -- MCU pipeline regression.} After migrating the bundled
Zephyr RTOS copy to a git submodule (Zephyr v3.7), the MIMaaS build pipeline
broke for two independent reasons: a missing kernel source file in the
TFLite-Micro CMake build (\texttt{strided\_slice\_common.cc}), and an upstream
TFLite-Micro change that dropped support for per-channel quantization in fully
connected layers -- a quantization scheme commonly produced by standard
TensorFlow Lite converters. Both had to be diagnosed and fixed without
breaking previously working models.

\textbf{Problem 2 -- Hailo-10H hardware backend and compatibility mapping.}
Unlike MCUs (which run an interpreter, TFLite-Micro, that executes arbitrary
graphs op-by-op), the Hailo-10H is a dataflow accelerator: a model must be
\emph{compiled ahead-of-time} by the Hailo Dataflow Compiler (DFC) into a
hardware-specific binary (\texttt{.hef}). This compilation involves graph
parsing, quantization/calibration, and static resource allocation across the
chip's compute clusters -- any of which can fail for a given architecture.
The problem was therefore not just ``build a backend that talks to the
board'', but also to answer the applied research question: \emph{for a
representative set of common CNN/Transformer/pose/detection architectures,
which ones can be compiled for Hailo-10H via DFC~5.3.0 out-of-the-box, and
what specifically causes the ones that fail to fail?}

\section{System Architecture}

\subsection{MIMaaS Core (existing, restructured)}

MIMaaS follows a producer/consumer architecture: a Flask REST API accepts
model uploads and enqueues a job; Redis acts as both the Celery message
broker and a live board-status store (heartbeat-based health tracking with a
30s TTL); and one Celery worker process per physical board consumes jobs from
a board-type queue and drives the actual hardware. During the internship, the
codebase was restructured from a flat script layout into an installable
Python package (\texttt{src/mimaas/}), with dedicated modules for the API,
worker, hardware backends, database models, and a YAML-based, environment-aware
configuration system (\texttt{config/default.yaml}, \texttt{development.yaml},
\texttt{production.yaml}).

\subsection{Hailo-10H Backend (new)}

The Hailo integration is implemented as a new hardware backend
(\texttt{src/mimaas/hardware/hailo\_utils.py}) invoked from the core pipeline
(\texttt{src/mimaas/core/pipeline.py}) whenever a request targets a board of
type \texttt{hailo\_rpi5}. Because the Hailo DFC toolchain requires a Linux
environment and is not natively supported on the Windows development machine,
the backend transparently shells out to WSL2, where DFC~5.3.0 is installed
inside a dedicated \texttt{pyenv} virtual environment. The pipeline for a
\texttt{.tflite} input is:

\begin{enumerate}
  \item \textbf{Parse}: \texttt{hailo parser tf model.tflite --hw-arch hailo10h}
    $\rightarrow$ \texttt{model.har} (Hailo Archive, an intermediate graph
    representation).
  \item \textbf{Optimize}: \texttt{hailo optimize model.har --use-random-calib-set}
    $\rightarrow$ \texttt{model\_opt.har} (quantization: mixed precision,
    statistics collection/calibration, hardware-parameter generation).
  \item \textbf{Compile}: \texttt{hailo compiler model\_opt.har --hw-arch hailo10h}
    $\rightarrow$ \texttt{model.hef} (resource partitioning across compute
    clusters/contexts, kernel compilation, final binary).
\end{enumerate}

If a \texttt{.hef} is supplied directly (e.g. a vendor-provided reference
model), the compile step is skipped entirely. Once a \texttt{.hef} exists, it
is transferred to the Raspberry Pi 5 over SSH/SCP (key-based, batch-mode
authentication) and executed with \texttt{hailortcli}, whose JSON/text output
(FPS, per-chip temperature) is parsed back into the same result schema used by
the MCU backends (status, inference time, RAM/ROM usage, power/energy where
available). This design means the REST API and result schema did not need to
change to support the new hardware class -- only the backend implementation.

\section{Methodology}

\subsection{MCU Track}

For the MCU track, the focus was on regression-fixing rather than new
feature work: (1) restoring the TFLite-Micro build after the Zephyr v3.7
submodule migration by patching the CMake kernel source list; (2) working
around the upstream removal of per-channel quantization support for dense
layers in TFLite-Micro, which is incompatible with MIMaaS's requirement to
accept arbitrary pre-converted \texttt{.tflite} uploads (users cannot be asked
to re-convert their models with
\texttt{\_experimental\_disable\_per\_channel\_quantization\_for\_dense\_layers});
and (3) fixing a metadata bug in the official \texttt{ppk2\_api} package that
prevented power measurement via the Nordic Power Profiler Kit~2 (PPK2).

\subsection{Hailo Compatibility Study}

For the Hailo track, a corpus of models was assembled from two sources:
(a)~14 vendor-provided, pre-compiled \texttt{.hef} reference models spanning
classification, object detection, and segmentation, used to validate the new
backend end-to-end on known-good binaries; and (b)~17 additional
classification architectures (MobileNet, ResNet, DenseNet, EfficientNet,
Inception/Xception, NASNet, ConvNeXt families) plus pose-estimation
(MoveNet) and audio (YAMNet) models, exported to \texttt{.tflite} from
Keras/TensorFlow and pushed through the full DFC pipeline described above.
Each model was tested sequentially (never in parallel, to avoid WSL2 memory
contention) and its outcome was classified as \emph{pass} (a \texttt{.hef} was
produced and executed on hardware, FPS/latency/temperature recorded) or
\emph{fail}, in which case the DFC error trace was inspected to attribute the
failure to one of the recurring root causes discussed in Section~5.2.

\section{Results}

\subsection{MCU Track}

All three MCU regressions were resolved with targeted, permanent fixes
committed to the Zephyr submodule / build configuration (documented in
\texttt{docs/WORKAROUNDS.md}), restoring build and flash compatibility for
the existing Nordic (nRF52833/nRF52840/nRF5340) and STM32 Nucleo boards
without requiring any change to already-deployed user models.

\subsection{Hailo Track}

\subsubsection{Reference models (pre-compiled \texttt{.hef})}

All 14 pre-compiled reference models executed successfully on the Hailo-10H
via the new backend (100\% pass rate), validating the SSH transfer,
\texttt{hailortcli} invocation, and result-parsing logic end-to-end.

\subsubsection{TFLite$\rightarrow$HEF compatibility matrix}

\begin{longtable}{@{}p{3.6cm}p{1.4cm}p{1.3cm}p{1.5cm}p{1.6cm}p{4.6cm}@{}}
\toprule
\textbf{Model} & \textbf{Status} & \textbf{FPS} & \textbf{Lat.\ (ms)} & \textbf{ROM (MB)} & \textbf{Notes} \\
\midrule
\endhead
MobileNetV1 & Pass & 1634.0 & 0.61 & 3.30 & Fastest model in the whole corpus \\
MobileNetV3-Small & Pass* & 466.6 & 2.14 & 3.27 & Fixed via 16-bit \texttt{avgpool} precision \\
MobileNetV3-Large & Pass & 301.2 & 3.32 & 5.76 & \\
ResNet50 & Pass & 261.8 & 3.82 & 19.86 & \\
ResNet50V2 & Pass & 239.2 & 4.18 & 19.82 & \\
DenseNet121 & Pass & 208.7 & 4.79 & 8.60 & \\
EfficientNetV2-B0 & Pass & 199.9 & 5.00 & 8.85 & \\
InceptionV3 & Pass & 166.6 & 6.00 & 19.50 & \\
DenseNet169 & Pass & 135.8 & 7.36 & 13.82 & \\
DenseNet201 & Pass & 105.7 & 9.46 & 18.85 & 43 min compile \\
Xception & Pass & 85.6 & 11.68 & 19.50 & Highest latency among passing models \\
InceptionResNetV2 & Pass & 68.1 & 14.68 & 39.78 & 2h17min compile (longest in corpus) \\
EfficientNetB0 & Pass* & 138.2 & 7.24 & 6.03 & Fixed via 16-bit \texttt{avgpool1} \\
EfficientNetB3 & \textbf{Fail} & -- & -- & -- & Same class as B0, \emph{not} fixable (Sec.~5.2.1) \\
NASNetMobile & \textbf{Fail} & -- & -- & -- & Non-scalar element-wise ratio (Sec.~5.2.2) \\
ConvNeXt-Tiny & \textbf{Fail} & -- & -- & -- & LayerNorm unsupported ops (Sec.~5.2.3) \\
YOLOv8n & \textbf{Fail} & -- & -- & -- & Exporter incompatibility (Sec.~5.2.4) \\
YOLOv8s & \textbf{Fail} & -- & -- & -- & Exporter incompatibility (Sec.~5.2.4) \\
MoveNet Lightning & \textbf{Fail} & -- & -- & -- & DFC translator bug (Sec.~5.2.5) \\
MoveNet Thunder & \textbf{Fail} & -- & -- & -- & DFC translator bug (Sec.~5.2.5) \\
YAMNet & \textbf{Fail} & -- & -- & -- & 1-D input unsupported \\
VGG16 & Abandoned & -- & -- & -- & WSL2 memory exhaustion, not a DFC limitation \\
\bottomrule
\caption{Compatibility results for TensorFlow/Keras models converted to TFLite
and compiled with Hailo DFC~5.3.0 for the Hailo-10H target. Entries marked
\emph{Pass*} required a model-script quantization override before compiling
successfully. All FPS/latency/ROM/temperature values were measured on
physical Hailo-10H hardware via \texttt{hailortcli}.}
\end{longtable}

Table~2 lists the 14 pre-compiled reference \texttt{.hef} models (100\% pass,
no DFC compilation involved), also measured on hardware, for completeness:

\begin{longtable}{@{}p{4cm}p{1.4cm}p{1.3cm}p{1.6cm}p{1.9cm}@{}}
\toprule
\textbf{Model} & \textbf{FPS} & \textbf{Lat.\ (ms)} & \textbf{HEF (MB)} & \textbf{Temp.\ (°C)} \\
\midrule
\endhead
cas\_vit\_t & 62.5 & 16.01 & 24.84 & 48.7 \\
deit\_small & 125.0 & 8.00 & 24.50 & 50.1 \\
hardnet39ds & 563.9 & 1.77 & 3.75 & 50.5 \\
hardnet68 & 219.0 & 4.57 & 16.63 & 51.3 \\
unet\_mobilenet\_v2 & 667.4 & 1.50 & 8.23 & 60.7 \\
deeplab\_v3\_mobilenet\_v2 & 50.6 & 19.78 & 2.52 & 56.9 \\
deeplab\_v3\_mobilenet\_v2\_wo\_dilation & 354.3 & 2.82 & 2.59 & 60.4 \\
fcn8\_resnet\_v1\_18 & 31.4 & 31.89 & 9.18 & 59.8 \\
stdc1 & 20.2 & 49.60 & 10.80 & 58.9 \\
ssd\_mobilenet\_v1 & 190.0 & 5.26 & 4.94 & 57.7 \\
tiny\_yolov4 & 562.0 & 1.78 & 4.54 & 64.8 \\
yolov5s\_bbox\_decoding\_only & 27.8 & 35.92 & 12.43 & 59.5 \\
yolo26m & 58.8 & 17.02 & 21.60 & 63.3 \\
yolov11x & 28.3 & 35.29 & 51.73 & 65.8 \\
\bottomrule
\caption{Reference \texttt{.hef} models used to validate the Hailo backend
end-to-end (classification, segmentation, and detection), all 100\% passing.}
\end{longtable}

Overall, of 17 newly-attempted architectures, 11 (65\%) compiled and executed
successfully out-of-the-box or with a small quantization override, while 6
(35\%) failed for reasons attributable to five distinct root causes.

\subsubsection{Novel-risk probes: targeted single-feature stress tests}

To close out the compatibility study, a final round of seven small,
purpose-built models was run, each isolating exactly \emph{one} architectural
feature not exercised by any model above (recurrence, native attention,
pre-quantized weights, pixel-shuffle upsampling, multi-input graphs, dynamic
batch dimensions, and transposed convolution via this study's own export
path, rather than a vendor pre-compiled \texttt{.hef}). Being deliberately
tiny (64--224\,px inputs), each either compiled and ran in under two minutes
or failed within seconds, making this an efficient way to map remaining
unknowns.

\begin{longtable}{@{}p{4cm}p{1.3cm}p{1.3cm}p{6.4cm}@{}}
\toprule
\textbf{Model} & \textbf{Status} & \textbf{FPS} & \textbf{Feature probed / outcome} \\
\midrule
\endhead
Tiny LSTM classifier & Pass & 3379.7 & Recurrent (LSTM) layer -- \emph{unexpectedly compiled}; with a fixed sequence length of 16 the TFLite converter most likely fully unrolled the recurrence into static ops, so this does not demonstrate true dynamic recurrence support \\
Tiny attention block & \textbf{Fail} & -- & Native \texttt{MultiHeadAttention} -- rejected on multiple simultaneous grounds: \texttt{BATCH\_MATMUL} unsupported (attention score/context matmuls), attention-reshape ops unresolvable, plus the same LayerNorm-via-moments issue as ConvNeXt (Sec.~5.2.3) \\
MobileNetV2 (INT8 pre-quantized) & \textbf{Fail} & -- & Fully-integer-quantized TFLite input -- explicit \texttt{UnsupportedQuantizedWeightsError}: DFC requires float32/16 weights and performs its own quantization; a model already quantized to int8 is rejected outright \\
Tiny ESPCN (super-resolution) & Pass & 867.4 & \texttt{DepthToSpace} (pixel-shuffle upsampling) -- compiled and ran without issue \\
Tiny Siamese (2 inputs) & Pass & 599.9 & Two independent graph inputs merged mid-graph -- supported, though compilation took $\sim$11\,min (longest of the seven, consistent with a more complex partitioning search) \\
MobileNetV2 (dynamic batch) & Pass & 1322.7 & Unspecified/dynamic leading batch dimension -- supported without modification \\
Tiny U-Net (\texttt{Conv2DTranspose}) & Pass & 2442.7 & Transposed convolution via this study's own TFLite$\rightarrow$DFC path (previously only seen in a vendor pre-compiled \texttt{.hef}) -- supported \\
\bottomrule
\caption{Novel-risk probes: seven small, single-feature models designed to
surface failure categories not covered by Sections 5.2.1--5.2.5. FPS/latency
measured on hardware via \texttt{hailortcli} for passing models.}
\end{longtable}

Two new, cleanly-attributable failure categories emerged: \textbf{(6) native
multi-head attention} (a combination of unsupported batched matrix
multiplication and the same LayerNorm decomposition issue already seen with
ConvNeXt) and \textbf{(7) pre-quantized INT8 input rejection} (DFC insists on
owning the quantization step itself and refuses an already-quantized model
outright, with a clear, actionable error message). Five of the seven probed
features -- multi-input graphs, dynamic batch dimensions, transposed
convolution, pixel-shuffle upsampling, and statically-unrollable recurrence
-- compiled without any workaround, indicating the DFC~5.3.0 parser's op
coverage is broader than the earlier classification-model failures alone
would suggest; failures are concentrated specifically around
\emph{quantization semantics} (avgpool shift-delta, non-scalar ratios,
pre-quantized weights) and \emph{LayerNorm/attention op decomposition},
rather than graph topology in general.

\section{Discussion}

\subsection{Root-cause taxonomy of DFC compilation failures}

\subsubsection{Global-average-pool quantization shift-delta overflow}
\textbf{Affected: MobileNetV3-Small (fixed), EfficientNetB0 (fixed),
EfficientNetB3 (not fixable).} The Hailo optimizer quantizes each
\texttt{avgpool} operation with an internal ``shift'' representing the scale
change between input and output; DFC enforces a maximum shift delta of~2.
Architectures whose global-average-pool operates over a wide dynamic range
(common in MobileNetV3's h-swish/SE blocks and EfficientNet's
squeeze-and-excitation blocks) can exceed this limit. For MobileNetV3-Small
(delta 2.87) and EfficientNetB0 (delta 4.73 on its single \texttt{avgpool1}),
forcing that specific layer to 16-bit precision via a model-script command
(\texttt{quantization\_param(avgpool1, precision\_mode=a16\_w16)}) resolved the
issue. For EfficientNetB3, which contains many more squeeze-and-excitation
blocks (each with its own average-pool), the same fix on \texttt{avgpool1}
immediately surfaced an identical error on \texttt{avgpool2} (delta~5.05);
patching every affected average-pool individually was judged impractical for
this study and the model was classified as non-fixable within the time
available. This suggests a general pattern: \emph{the more squeeze-excitation
blocks a model has, the less likely a single-layer quantization override is
sufficient}, and a systematic script (auto-detecting and patching all
average-pool layers) would be needed for larger EfficientNet variants.

\subsubsection{Non-scalar element-wise-multiply ratio}
\textbf{Affected: NASNetMobile.} A distinct quantization failure in an
\texttt{ElementWiseMult} layer, where DFC's optimizer assumes the ratio
between output and input scale is a scalar, but NASNetMobile's branch
structure produces a per-channel ratio. Unlike the average-pool issue above,
forcing 16-bit precision did not resolve this failure, indicating an
architectural incompatibility rather than a precision-budget problem.

\subsubsection{Unsupported LayerNorm decomposition (ConvNeXt family)}
\textbf{Affected: ConvNeXt-Tiny.} ConvNeXt's LayerNorm is exported by the
TFLite converter as an explicit \texttt{moments} computation (mean +
\texttt{SquaredDifference} + variance), rather than a fused normalization op.
The Hailo parser rejects both the \texttt{SQUARED\_DIFFERENCE} op (unsupported
entirely) and any \texttt{ReduceMean} whose reduction axis is the channel axis
(DFC requires reductions to be spatial, i.e.\ equivalent to a global average
pool). This is a graph-\emph{translation} failure at the parser stage
(occurring within 20 seconds, before any quantization work begins), not a
numerical/quantization failure, and is not fixable without re-exporting the
model with a different LayerNorm implementation or waiting for parser support.
This is classified as \textbf{non-fixable} in the current DFC version.

\subsubsection{Exporter-introduced unsupported ops (YOLOv8)}
\textbf{Affected: YOLOv8n, YOLOv8s.} Both failed during parsing with
\texttt{UnsupportedShuffleLayerError} / \texttt{UnsupportedModelError},
traced to the Ultralytics \texttt{litert\_torch} TFLite exporter producing a
detection head with operations (\texttt{TRANSPOSE}, \texttt{DFL}, \texttt{EWAdd})
and a channel-shape mismatch inside a C2f block that the parser cannot
translate. This is an \emph{export-tool} incompatibility rather than a DFC or
model-architecture limitation; a different export path (e.g.\ ONNX
$\rightarrow$ Hailo's dedicated YOLO parsing flow, which several published
Hailo Model Zoo YOLO variants use) would likely avoid this failure mode
entirely.

\subsubsection{Internal DFC translator bug on non-4D tensors (MoveNet)}
\textbf{Affected: MoveNet Lightning, MoveNet Thunder.} Both crash with an
\texttt{IndexError} deep inside Hailo's own \texttt{format\_conversion.py},
reproduced twice with identical stack traces and unaffected by toggling
\texttt{SELECT\_TF\_OPS} at export time. This points to a bug in Hailo's
internal format-conversion layer when handling MoveNet's non-4D intermediate
tensor shapes, rather than an issue in the exported model itself; it is
plausible this is fixed in a newer DFC release, but this was not verifiable
within DFC~5.3.0.

\subsubsection{Native multi-head attention (transformer blocks)}
\textbf{Affected: a minimal \texttt{MultiHeadAttention} probe model.} Beyond
the LayerNorm-via-moments issue already characterized for ConvNeXt, attention
blocks fail for an additional, independent reason: the batched matrix
multiplications used to compute attention scores and context vectors
(\texttt{BATCH\_MATMUL}) are unsupported outright by the DFC~5.3.0 parser, as
are several of the reshape/einsum operations Keras's \texttt{MultiHeadAttention}
emits internally. This is therefore a \textbf{compound failure}: even if the
LayerNorm issue were worked around, native attention would still fail on the
matmul/reshape ops. Transformer-style architectures should be considered
unsupported by this DFC version until at least one of these two independent
issues is resolved upstream.

\subsubsection{Pre-quantized (INT8) TFLite input rejection}
\textbf{Affected: a MobileNetV2 export using full-integer post-training
quantization.} DFC is designed to own the quantization step itself (via its
Optimizer stage) and explicitly rejects a \texttt{.tflite} file whose weights
are already quantized to int8/int32, raising a clear
\texttt{UnsupportedQuantizedWeightsError} at the parser stage rather than
attempting a (likely incorrect) double quantization. This is not a bug but a
deliberate design constraint: \textbf{MIMaaS users targeting Hailo-10H must
supply float32/float16 \texttt{.tflite} exports}, never a pre-quantized
INT8 model, which is a useful, actionable constraint to surface to end users
of a future Hailo-enabled MIMaaS endpoint.

\subsection{Environment constraints}

Two failures were not architectural at all: \textbf{YAMNet} was rejected
because DFC requires a 4D (image-like) input tensor and YAMNet's native input
is a 1D waveform; and \textbf{VGG16} was abandoned after $>$1h40min of
compilation without completion, attributable to the interaction of its
138M-parameter, dense-layer-heavy architecture (a 103M-parameter first dense
layer alone) with the shared, memory-constrained WSL2 VM (14GB available),
rather than a hard DFC resource limit -- large dense networks should be
retried on a machine with more available memory or a native Linux
installation of DFC before being written off entirely.

\subsection{Compilation time variability}

Compilation time varied by more than two orders of magnitude across the
corpus tested: from single-digit minutes for smaller MobileNet/DenseNet
variants up to \textbf{2 hours 17 minutes} for InceptionResNetV2, driven
primarily by the multi-context partition search (DFC's iterative search for
a valid layer-to-cluster assignment across multiple hardware execution
contexts). This has direct practical implications for how such a
compatibility study -- or a future automated MIMaaS Hailo endpoint -- should
be scheduled: compilation cannot be assumed to complete within a fixed
timeout, and strictly sequential (non-parallel) execution is necessary on a
single shared WSL2 VM to avoid out-of-memory failures during concurrent
compilations.

\section{Conclusion and Future Work}

This internship delivered two concrete outcomes. First, the existing
MIMaaS MCU pipeline was repaired and restructured into a maintainable Python
package, resolving a post-migration regression that had left three board
families unable to build or run TFLite-Micro models. Second, MIMaaS was
successfully extended to a new class of hardware -- the Hailo-10H edge NPU --
via a new backend that reuses the platform's existing API/result contract,
and this new backend was used to run the first systematic Hailo-10H DFC~5.3.0
compatibility study across 30+ well-known architectures, producing (a)~a
reference performance table for models that compile out-of-the-box and
(b)~a root-cause taxonomy for the ones that do not, distinguishing
\emph{fixable} quantization issues (single average-pool overrides),
\emph{architecturally non-fixable} issues (non-scalar ratios, unsupported
LayerNorm decomposition, native attention), \emph{exporter-specific} issues
(YOLOv8), \emph{compiler-internal bugs} (MoveNet), and \emph{deliberate
design constraints} (rejection of pre-quantized INT8 inputs).

Recommended future work: (1)~automate detection and patching of \emph{all}
problematic average-pool/SE-block layers in a model script, rather than a
single hard-coded layer name, to recover larger EfficientNet variants;
(2)~evaluate an ONNX-based export path for YOLO-family detectors, which is
better supported by Hailo's tooling than the current TFLite exporter;
(3)~re-run the MoveNet and VGG16 cases against a newer DFC release and/or a
native Linux host to separate compiler bugs and environment limits from true
architectural incompatibilities; and (4)~expose the Hailo backend through the
same MIMaaS REST API used for MCU boards, so that end users can request
Hailo-10H evaluations without needing direct WSL2/DFC access, using the
compatibility matrix in Section~5 to set user expectations up front.

\section{Personal Reflection}

\textit{[Draft below -- please edit freely, this is meant as a starting
point reflecting what was actually done, not a template to fill mechanically.]}

Coming into this internship with some general programming background but no
prior hands-on exposure to embedded ML deployment, edge NPUs, or production-grade
service architecture, three things stood out as the most valuable learning
outcomes.

\textbf{Reading a system before changing it.} The first real task was not
writing new code but understanding and repairing an existing pipeline
(MIMaaS's MCU backend) after a dependency migration broke it. This required
tracing a build error and a runtime error back through several layers --
CMake source lists, a third-party TFLite-Micro upstream change, a hardware
API library -- before a one-line fix could be applied with confidence. The
lesson that stuck: most real engineering work is understanding \emph{why}
something that used to work no longer does, not writing code from a blank
page.

\textbf{Ahead-of-time compilation is a fundamentally different failure mode
than interpreted inference.} Prior experience with TensorFlow Lite had been
limited to running an interpreter, where a model either loads or does not,
and errors are usually about missing/unsupported ops. Working with the Hailo
Dataflow Compiler was a different experience entirely: a model can pass
parsing, pass calibration, and still fail deep inside \emph{quantization
arithmetic} (the average-pool shift-delta errors) or during \emph{static
resource allocation} (multi-context partition search taking over an hour for
some architectures). This distinction -- between ``does the graph contain
supported operations'' and ``can this graph be statically compiled onto a
fixed amount of hardware'' -- is, in hindsight, the core conceptual
difference between deploying to a general-purpose interpreter (MCU/CPU) and
deploying to a dataflow accelerator (NPU), and it shaped how the whole
compatibility study was organized (grouping failures by \emph{which stage of
the toolchain} they occurred in, rather than just pass/fail).

\textbf{Negative results are still results, if the cause is precise enough.}
Roughly a third of the new architectures tested failed to compile. Early on,
the instinct was to see this as simply ``incomplete work''. Reframing each
failure as a small root-cause investigation (attributing NASNetMobile,
ConvNeXt, YOLOv8, and MoveNet to four \emph{different} causes, rather than a
single ``Hailo doesn't support X'' bucket) turned out to be far more useful,
both for this report and, hopefully, for anyone at AnkiLab who later tries
one of these architectures again on a newer DFC version: knowing precisely
\emph{why} something fails tells you whether it is worth retrying later
(a compiler bug, possibly fixed upstream), worth fixing on the export side
(YOLOv8's exporter), or a genuine current limitation of the hardware/toolchain
(non-scalar quantization ratios, unsupported LayerNorm decomposition).

More generally, this internship was a good introduction to the gap between
``a model works on my laptop'' and ``a model is deployable on a specific
piece of edge hardware'' -- a gap that, in industry, is exactly what a
platform like MIMaaS exists to make legible to non-hardware-expert users.

\appendix
\section*{Appendix: Reproducing the Results}
See the accompanying \emph{MIMaaS + Hailo-10H Setup Manual} for full
installation and reproduction instructions. All quantitative results in
Section~5 were produced via the \texttt{flash\_and\_evaluate()} entry point in
\texttt{src/mimaas/core/pipeline.py}, which is hardware-agnostic and dispatches
to \texttt{src/mimaas/hardware/hailo\_utils.py} for boards of type
\texttt{hailo\_rpi5}.

\end{document}
